\documentclass[12pt,a4paper]{article}
\usepackage[utf8]{inputenc}
\usepackage[vietnamese]{babel}
\usepackage[left=2cm,right=2cm,top=2cm,bottom=2cm]{geometry}
\usepackage{mathptmx}
\usepackage{enumitem}

\pagestyle{plain}
\linespread{1.15}

\begin{document}

% --- HEADER QUỐC HIỆU & CƠ QUAN BAN HÀNH ---
\noindent
\begin{minipage}[t]{0.45\textwidth}
    \begin{center}
        \small
        CÔNG AN THÀNH PHỐ HÀ NỘI\\
        \textbf{\underline{PHÒNG CẢNH SÁT HÌNH SỰ}}\\[0.1cm]
        Số: 03/BC-CSHS\\
        \textit{V/v tiến độ điều tra ban đầu vụ án Khang}
    \end{center}
\end{minipage}
\hfill
\begin{minipage}[t]{0.52\textwidth}
    \begin{center}
        \small
        \textbf{CỘNG HÒA XÃ HỘI CHỦ NGHĨA VIỆT NAM}\\
        \textbf{\underline{Độc lập - Tự do - Hạnh phúc}}\\[0.1cm]
        \textit{Hà Nội, ngày 25 tháng 07 năm 2026}
    \end{center}
\end{minipage}

\vspace{0.6cm}

% --- TIÊU ĐỀ VĂN BẢN ---
\begin{center}
    \textbf{\large BÁO CÁO TIẾN ĐỘ ĐIỀU TRA BAN ĐẦU}\\
    \textit{(Chuyên án TEST-99 — Án mạng tại Phân khu Cảng)}
\end{center}

\vspace{0.3cm}

\noindent Ban Chuyên án báo cáo Thủ trưởng Cơ quan CSĐT Công an TP. Hà Nội tổng hợp tiến độ triển khai lực lượng và kết quả lấy lời khai các đối tượng liên quan vụ án mạng xảy ra ngày 24/07/2026 tại số 14 Đường Bờ Sông như sau:

\section*{I. KẾT QUẢ RÀ SOÁT DẤU VẾT VẬT CHỨNG \& TRÍCH XUẤT THÔNG TIN}
\begin{enumerate}[label=\arabic*., leftmargin=1.5cm]
    \item \textbf{Khám nghiệm hiện trường \& Tử thi:} Nạn nhân Khang bị tổn thương va đập chẩm gáy lúc khoảng 20:00 và bị vết đâm trực tiếp vào động mạch cảnh gây mất máu cấp tử vong muộn hơn.
    \item \textbf{Trích xuất tin nhắn điện thoại \& Vật chứng:} Phát hiện tin nhắn trao đổi gay gắt hẹn gặp tối 24/07 giữa Khang và Mai về chuyện di chúc đền bù, cùng chuỗi tin nhắn khất nợ dồn dập từ số di động rác chưa lưu tên \texttt{0985.xxx.341}. Tại hiện trường thu giữ 01 hóa đơn nạp tiền di động vò nát cho số máy này dưới gầm bàn.
\end{enumerate}

\section*{II. TỔNG HỢP LỜI KHAI CỦA NGUYÊN ĐƠN \& CÁC ĐỐI TƯỢNG LIÊN QUAN}
\begin{enumerate}[label=\arabic*., leftmargin=1.5cm]
    \item \textbf{Nguyễn Thị Lụa (Hàng xóm sát nhà):} Khai thấy Mai nổ máy xe về một mình lúc 19:00, Vũ đi bộ ra đầu ngõ; nghe tiếng vỡ đồ đạc phát ra từ nhà Khang lúc hơn 20:15 đêm 24/07.
    \item \textbf{Trần Thị Hà (Bạn gái nạn nhân):} Khai ở nhà xem tivi một mình cả tối. Khai tình cảm hai người đằm thắm và sực nhớ cung cấp thông tin nghi vấn về Nguyễn Thanh Tùng dạo này lởn vởn quanh ngõ.
    \item \textbf{Trần Ngọc Mai (Em họ):} Khai sang nhà Khang lúc 18:30 bàn chuyện di chúc đền bù rồi đi xe máy về một mình lúc 19:00 (vợ chồng đi chung xe, Vũ bảo có bạn đón đi nhậu).
    \item \textbf{Lê Quang Vũ (Chồng Mai):} Khai bảo Mai lấy xe máy về trước lúc 19:00 để đứng chờ bạn đi ô tô đón ra quán bia ven sông ngồi uống từ 19:40 đến 21:30, không có người làm chứng.
    \item \textbf{Nguyễn Thanh Tùng (Bạn từ nhỏ):} Ban đầu chìa cuống vé xe khách 19:30 để chối tội. Khi bị vạch trần thì lúng túng thừa nhận sang nhà Khang lúc 20:00 đòi lại đồ cũ nhưng chối không hại Khang.
\end{enumerate}

\section*{III. KẾT QUẢ XÁC MINH NHÁNH NGHI VẤN PHỤ}
Đội Cảnh sát hình sự đã tiến hành triệu tập và lấy lời khai xác minh 02 đối tượng có hiềm khích phụ do nhân chứng Trần Thị Hà cung cấp:
\begin{enumerate}[label=\arabic*., leftmargin=1.5cm]
    \item \textbf{Nguyễn Văn Hùng (Sinh năm 1968, Trú tại ngõ sau nhà nạn nhân):} Mâu thuẫn tranh chấp rào đất đền bù. Xác minh từ 19:00 đến 21:30 ngày 24/07/2026 tham dự cuộc họp Tổ dân phố số 4 tại Nhà văn hóa (có biên bản và hơn 30 hộ dân làm chứng) $\rightarrow$ \textbf{Bằng chứng ngoại phạm hợp lệ, rõ ràng; đã chính thức loại trừ nghi vấn.}
    \item \textbf{Nguyễn Hoài Nam (Sinh năm 1993, Trú tại Phân khu Cảng):} Mâu thuẫn nợ nần tài chính. Trích xuất camera Clb Billiards X-Club xác định đối tượng có mặt thi đấu liên tục từ 18:30 đến 23:00 ngày 24/07/2026 $\rightarrow$ \textbf{Bằng chứng ngoại phạm hợp lệ, rõ ràng; đã chính thức loại trừ nghi vấn.}
\end{enumerate}

\section*{IV. ĐÁNH GIÁ VÀ HƯỚNG ĐIỀU TRA TIẾP THEO}
Tất cả các đối tượng chính (Mai, Vũ, Tùng) đều không có bằng chứng ngoại phạm hợp lệ và đều che giấu động cơ thực sự. Ban Chuyên án tiếp tục xác minh chi tiết thời gian biểu và giám định dấu vết sinh học tại hiện trường.

\vspace{0.8cm}

\noindent
\begin{minipage}[t]{0.45\textwidth}
    \small
    \textbf{\textit{Nơi nhận:}}\\
    - Thủ trưởng Cơ quan CSĐT (b/c);\\
    - Viện KSMND TP. Hà Nội;\\
    - Trưởng Ban Chuyên án;\\
    - Lưu: VT, CSHS.
\end{minipage}
\hfill
\begin{minipage}[t]{0.5\textwidth}
    \begin{center}
        \small
        \textbf{TRƯỞNG BAN CHUYÊN ÁN}\\
        \textit{(Ký, ghi rõ họ tên)}\\[1.5cm]
        \textbf{Đại úy Lê Minh}
    \end{center}
\end{minipage}

\end{document}
